% !TeX root = ../papers/nc-p3.tex
\section{十五次元閉包の全正則整域窓における四次中心パラメータ}
\label{sec:main-domain-quartic-parameters}
\begin{lemma}\label{lem:main-center-integral-degrees}
十五次元閉包の有限四辺の正則整域窓で、射影中心の各整曲線 $C$ に対し $HC$ は正整数、
各整曲面 $V$ に対し $H^2 V$ は正整数である。さらに固定成分のない可動な因子鉛筆の一般の
有効因子 $V$ について $H^2V\ge2$ である。
\end{lemma}
\begin{proof}
一般の安定化はどの曲線でも自明である（定理\ref{thm:octic-surface-isolated-inertia}）。
曲線の体 $K=k(C)$ 上で中心指標の次数を有理的に自明化する。
特殊化された被約traceの次元ベクトル $4\delta$ の半単純表現を $K$ 上で取れる。
実際代数閉拡大後には全単純の正整数重複つきの和である（中心有限なtrace環の格子による特殊化）。
重複度は特殊化traceから一意でGalois不変である。有限次元繊維の半単純商の各中心は $K$ の有限拡大で、
その上の行列環に分裂する（Tsenの定理）。従って各軌道の絶対単純の和が $K$ 上で定義され、指定の重複度で取れる。
曲面 $V$ の一般点はAzumayaであり、体 $k(V)$ での指数を $r\mid4$ とすれば
一つの単純の頂点次元は $4r$ である。

これらをそれぞれ台の上の連接な窓加群に延ばす。詳しくは、一般点の中心次数を自明化した
開集合で行列を広げ、次数作用に対応する束として延長し、有限な代数の作用で閉じる。
有限加群性により連接である。次数安定化の点も含めるときは斉次錐上で各有理生成元に分母を掛け、
得た有限個の斉次元で生成する部分加群を取ってよい。これは
定理\ref{thm:central-cut-division-index}の一般単純の延長と同じで、台を $C,V$ に置き換えたものである。
高い次数を取れば各剰余位置での偶数中心移動は $2H$ による捩りになる。
有限生成頂点射影分解により、整数添字に沿った一つのHilbert級数は $p(t)/(1-t)^4$、
$p$ は整数係数Laurent多項式である。曲線なら一次成長であるから $(1-t)^2$ を除いた後の
最高係数は正整数 $e$、すなわち経路長 $8n+i$ で $8en+O(1)$ である。
同じ係数を曲線上で取ると $4(2H\cdot C)n$ であり $HC=e$ となる。
曲面なら $64e n^2/2$ と $4r(2H)^2V n^2/2$ が一致し、$H^2V=4e/r$ である。

可動鉛筆で $H^2V=1$ と仮定する。定理\ref{thm:finite-center-divisor-torsion}より
$V\sim_{\mathbb Q}H/4$。独立に選ぶ幾何学的一般 $S\in|2H|$ の上で、この制限のWeil因子は
$h/4$ に有理線型同値である（滑らかな部分で制限して延長する）。
最小解消で一般のこの曲線の狭義変換を $G$、例外 $(-2)$ 曲線を $E_j$ とする。
一般切断は鉛筆の固定台のどの曲線も含まないため、狭義変換の鉛筆にも固定曲線はなく $G^2\ge0$。
二つの一般メンバーを交叉させればよい。各例外への交点数は一般メンバーで一定である。
$A_1$ の直交引き戻しを用いると
\[
             0\le G^2=\tfrac12-\tfrac12\sum_j(GE_j)^2 .
\]
左辺は整数なので、ただ一つの $j$ でだけ $GE_j=1$ となる。しかしこの個数は1になれない。
鉛筆の一般メンバーを先に定めてその定義体を代数閉にし、その後で $S$ の独立な一般体を取る。
切断の特異点は $X$ の特異曲線との横断点に限られる。各整特異曲線上の普遍超平面接続は
射影空間束で整なので、この曲線から来る点は一般体上で一つのGalois軌道をなす。
その大きさは $2H\cdot C$ で、先の整除により少なくとも2である。
最小解消と狭義変換の交点数は共役で不変なので、唯一の例外点の存在は矛盾する。
\end{proof}

\begin{theorem}\label{thm:main-domain-quartic-parameters}
十五次元閉包の有限四辺で、\emph{全ての}正則整域窓の四次中心系は基点無しである。
従って実現中心は通常の四次超曲面または整な三次元二次曲面の四次分岐二重被覆で、
偏極中心と零番Azumaya隅のKummer印が添字保持同型の完全不変量である。
この閉包内で正則整域性は経路長4の四中心元を使う有限検査と同値であり、
経路次元は14次、中心剰余は13次、標準複体は16次までで足りる。
\end{theorem}
\begin{proof}
まず三次元二次曲面上の基点付き表示 $q=0,\ w\ell+f=0$ は階数を問わず不可能である。
射影四空間で $q=\ell=f=0$ の点 $y$ がある（三式なので空にならない）。
中心に $[t:w]\mapsto[t y:w]$ という重み $(1,2)$ の射影直線が閉埋込みされる。
実際 $y$ の一つの非零座標から $t$ が戻り、斉次座標環の射は全射である。
この曲線上で $\mathcal O(2)$ の次数は1（$[t^2:w]$ が普通の射影直線との同型）で、補題に反する。
これは前節の分岐による必要条件を満たす候補も除く。

曲面像のうち四次は定理\ref{thm:exclude-del-pezzo-center-images}で除いた。
三次スクロールなら一般の罫線 $R_t$ の逆像の因子部分 $V_t$ は非零な可動鉛筆である。
実際スクロールの頂点があればそれを除いて罫線パラメータへの写像を合成し、
その一般の繊維を閉じればよい。各 $V_t$ は一般の罫線を支配し、一般メンバーに共通な因子は除いている。
よって $H^2V_t\ge2$。錐型では頂点を通る一般超平面は異なる三罫線の和であり、
四次中心切断の因子にはこれらの一般繊維が全て入る。従って $4=H^3\ge 3\cdot2$ となり不可能である。

滑らかな三次スクロールでは負の切断である直線を $N$ とすれば超平面因子類は $N+2R$ である。
ここで $N$ が四次有理写像の定義域から抜け落ちないことを示す。
四次系の一次元基底台はただ一本の整曲線 $B_0$ で、一般点で被約である。
実際一般八次曲面における基底スキームは異なる単純二点なので、その一次元成分の
重複度つき $2H$ 次数の和は2である。各整曲線の $2H$ 次数は2以上であった。
一般基点は滑らかな $X$ にある。非Azumaya基点は次数安定化点であって孤立するからである。
一般 $p\in B_0$ で五切断の一次法線方向の値は二次元を張り、一つの直線 $L_p\subset Y$ を定める。
一般の $S$ が $B_0$ と横断するところで、単純基点を一回吹き上げた例外がこの直線へ同型に写るためである。
一次法線写像なので $L_p$ は $p$ だけで決まる。

$S$ の二基点 $P,Q$ から得る二直線は異なる。
事実二点を吹き上げ解消した標準写像は次数2であり、どちらの例外直線でもその一般点で分岐している。
一般超平面で横断すると、例外直線は像の直線と同型で、その交点は前掲の一般標準曲線上のWeierstrass点である。
もし曲面の写像がそこでエタールなら曲線の制限もエタールになる。
標準写像のStein分解では各例外直線はその像の因子を支配し、分岐指数2の因子が二本とも
同じ因子上には来られない。よって二直線は異なる。
従って整な $B_0$ の一般点で $L_p$ が固定した $N$ であることはなく、一般の $S$ ではどちらも $N$ でない。
解消した $S$ は $Y$ へ固有全射であるから、$N$ の一般点の逆像は基点を除いた写像の定義域にある。
元の $A_1$ の解消例外は $S$ の標準写像で点にしか写らないことも使った。
従って $X$ の定義域におけるCartier因子 $N$ の引き戻しを閉じると非零な有効因子を得る。
$N$ と一般の二罫線の和に対応する超平面を選ぶと $H^3\ge1+2\cdot2$ となって矛盾する。
これで全ての基点付き候補が除かれた。

四つの一般四次元がパラメータになるから、既出の中心型と
定理\ref{thm:central-cut-kummer-classification}が例外も含め適用できる。
逆に四次パラメータ検査が通れば正則かつ中心有限で合成整域である。
最終の三上限は命題\ref{prop:central-arrow-module-cutoff}で $L=12$ として得る。
\end{proof}
